\documentclass[12pt]{article}

\usepackage[T1]{fontenc}
\usepackage{lmodern}
\usepackage{amsmath,amssymb,mathtools}
\usepackage[margin=1in]{geometry}
\usepackage{microtype}

\setlength{\parindent}{0pt}
\setlength{\parskip}{0.6em}
\setlength{\abovedisplayskip}{0.9em}
\setlength{\belowdisplayskip}{0.9em}

\newcommand{\R}{\mathbb{R}}
\newcommand{\e}{\mathrm{e}}
\newcommand{\bigO}{\mathcal{O}}

\title{Linearization of the prediction equation}
\author{}
\date{}

\begin{document}
\maketitle

\begin{equation}
    \dot{x}(t)=A_{\sigma(t)}x(t)+b_{\sigma(t)},
    \qquad x(t)\in\R^n.
    \label{eq:affine-system}
\end{equation}

\begin{equation}
    \chi(t)=
    \begin{bmatrix}
        x(t)\\
        1
    \end{bmatrix}.
    \label{eq:augmented-state}
\end{equation}

\begin{equation}
\begin{aligned}
    \dot{\chi}(t)
    &=
    \begin{bmatrix}
        \dot{x}(t)\\[2pt]
        \dfrac{\mathrm d}{\mathrm dt}(1)
    \end{bmatrix}                                                    \\[4pt]
    &=
    \begin{bmatrix}
        A_{\sigma(t)}x(t)+b_{\sigma(t)}\\
        0
    \end{bmatrix}                                                    \\[4pt]
    &=
    \begin{bmatrix}
        A_{\sigma(t)} & b_{\sigma(t)}\\
        0_{1\times n} & 0
    \end{bmatrix}
    \begin{bmatrix}
        x(t)\\
        1
    \end{bmatrix}.
\end{aligned}
\label{eq:augmented-construction}
\end{equation}

\begin{equation}
    \widetilde A_{\sigma(t)}=
    \begin{bmatrix}
        A_{\sigma(t)} & b_{\sigma(t)}\\
        0_{1\times n} & 0
    \end{bmatrix}.
    \label{eq:augmented-matrix}
\end{equation}

\begin{equation}
    \boxed{
        \dot{\chi}(t)=\widetilde A_{\sigma(t)}\chi(t)
    }.
    \label{eq:augmented-system}
\end{equation}

\begin{equation}
    0=t_0<t_1<t_2<\cdots<t_N,
    \label{eq:boundaries}
\end{equation}

\begin{equation}
    \sigma(t)=q_i,
    \qquad t_{i-1}\leq t<t_i,
    \qquad i=1,2,\ldots,N.
    \label{eq:mode-intervals}
\end{equation}

\begin{equation}
    F_i=\widetilde A_{q_i},
    \qquad i=1,2,\ldots,N.
    \label{eq:Fi}
\end{equation}

\begin{equation}
    \Delta t_i=t_i-t_{i-1}.
    \label{eq:actual-duration}
\end{equation}

\begin{equation}
    \chi(t_i)=\e^{F_i\Delta t_i}\chi(t_{i-1}).
    \label{eq:one-interval-propagation}
\end{equation}

\begin{equation}
\boxed{
    \chi(t_N)
    =\e^{F_N\Delta t_N}
     \e^{F_{N-1}\Delta t_{N-1}}
     \cdots
     \e^{F_2\Delta t_2}
     \e^{F_1\Delta t_1}
     \chi(t_0)
    }.
    \label{eq:actual-cycle}
\end{equation}

\begin{equation}
    \Delta\bar t_i=\bar t_i-\bar t_{i-1},
    \qquad i=1,2,\ldots,N.
    \label{eq:nominal-duration}
\end{equation}

\begin{equation}
    \Delta t_i=\Delta\bar t_i+\delta t_i.
    \label{eq:duration-perturbation}
\end{equation}

\begin{equation}
\boxed{
    \bar\chi(t_N)
    =\e^{F_N\Delta\bar t_N}
     \e^{F_{N-1}\Delta\bar t_{N-1}}
     \cdots
     \e^{F_2\Delta\bar t_2}
     \e^{F_1\Delta\bar t_1}
     \bar\chi(t_0)
    }.
    \label{eq:nominal-cycle}
\end{equation}

\begin{equation}
    e(t)=\chi(t)-\bar\chi(t).
    \label{eq:error}
\end{equation}

\begin{equation}
    \chi(t_0)=\bar\chi(t_0)+e(t_0).
    \label{eq:initial-state-decomposition}
\end{equation}

\begin{equation}
\begin{aligned}
    e(t_N)
    ={}&
    \e^{F_N(\Delta\bar t_N+\delta t_N)}
    \e^{F_{N-1}(\Delta\bar t_{N-1}+\delta t_{N-1})}
    \cdots                                                        \\
    &\quad
    \e^{F_2(\Delta\bar t_2+\delta t_2)}
    \e^{F_1(\Delta\bar t_1+\delta t_1)}
    \left(\bar\chi(t_0)+e(t_0)\right)                            \\
    &-
    \e^{F_N\Delta\bar t_N}
    \e^{F_{N-1}\Delta\bar t_{N-1}}
    \cdots
    \e^{F_2\Delta\bar t_2}
    \e^{F_1\Delta\bar t_1}
    \bar\chi(t_0).
\end{aligned}
\label{eq:exact-error}
\end{equation}

\begin{equation}
    \frac{\mathrm d}{\mathrm ds}\e^{F_i s}
    =F_i\e^{F_i s}.
    \label{eq:exponential-derivative}
\end{equation}

\begin{equation}
\begin{aligned}
    \e^{F_i(\Delta\bar t_i+\delta t_i)}
    ={}&\e^{F_i\Delta\bar t_i}
       +F_i\e^{F_i\Delta\bar t_i}\,\delta t_i
       +\bigO(\delta t_i^2).
\end{aligned}
\label{eq:single-exponential-expansion}
\end{equation}

\begin{equation}
    \phi_i=\e^{F_i\Delta\bar t_i}.
    \label{eq:phi}
\end{equation}

\begin{equation}
    \e^{F_i(\Delta\bar t_i+\delta t_i)}
    =\phi_i+F_i\phi_i\,\delta t_i+\bigO(\delta t_i^2).
    \label{eq:single-exponential-phi}
\end{equation}

\begin{equation}
\begin{aligned}
&\e^{F_N(\Delta\bar t_N+\delta t_N)}
 \e^{F_{N-1}(\Delta\bar t_{N-1}+\delta t_{N-1})}
 \cdots
 \e^{F_1(\Delta\bar t_1+\delta t_1)}                              \\
&\quad\doteq
 \left(\phi_N+F_N\phi_N\,\delta t_N\right)
 \left(\phi_{N-1}+F_{N-1}\phi_{N-1}\,\delta t_{N-1}\right)
 \cdots
 \left(\phi_1+F_1\phi_1\,\delta t_1\right).
\end{aligned}
\label{eq:product-before-expansion}
\end{equation}

\begin{equation}
\begin{aligned}
&\left(\phi_N+F_N\phi_N\,\delta t_N\right)
 \left(\phi_{N-1}+F_{N-1}\phi_{N-1}\,\delta t_{N-1}\right)
 \cdots
 \left(\phi_1+F_1\phi_1\,\delta t_1\right)                       \\
&\quad\doteq
 \phi_N\phi_{N-1}\cdots\phi_2\phi_1                            \\
&\qquad
 +F_N\phi_N\phi_{N-1}\cdots\phi_2\phi_1\,\delta t_N          \\
&\qquad
 +\phi_NF_{N-1}\phi_{N-1}\cdots\phi_2\phi_1\,\delta t_{N-1}  \\
&\qquad
 +\phi_N\phi_{N-1}F_{N-2}\phi_{N-2}\cdots
  \phi_2\phi_1\,\delta t_{N-2}                                 \\
&\qquad
 +\cdots                                                         \\
&\qquad
 +\phi_N\phi_{N-1}\cdots\phi_3F_2\phi_2\phi_1\,\delta t_2   \\
&\qquad
 +\phi_N\phi_{N-1}\cdots\phi_3\phi_2F_1\phi_1\,\delta t_1.
\end{aligned}
\label{eq:ordered-product-expansion}
\end{equation}

\begin{equation}
\begin{aligned}
    e(t_N)
    \doteq{}&
    \phi_N\phi_{N-1}\cdots\phi_2\phi_1
    \left(\bar\chi(t_0)+e(t_0)\right)                            \\
    &+F_N\phi_N\phi_{N-1}\cdots\phi_2\phi_1
      \left(\bar\chi(t_0)+e(t_0)\right)\delta t_N                \\
    &+\phi_NF_{N-1}\phi_{N-1}\cdots\phi_2\phi_1
      \left(\bar\chi(t_0)+e(t_0)\right)\delta t_{N-1}          \\
    &+\cdots                                                     \\
    &+\phi_N\phi_{N-1}\cdots\phi_3\phi_2F_1\phi_1
      \left(\bar\chi(t_0)+e(t_0)\right)\delta t_1              \\
    &-\phi_N\phi_{N-1}\cdots\phi_2\phi_1\bar\chi(t_0).
\end{aligned}
\label{eq:error-before-cancellation}
\end{equation}

\begin{equation}
\begin{aligned}
    e(t_N)
    \doteq{}&
    \phi_N\phi_{N-1}\cdots\phi_2\phi_1e(t_0)                   \\
    &+F_N\phi_N\phi_{N-1}\cdots\phi_2\phi_1
      \bar\chi(t_0)\,\delta t_N                                 \\
    &+\phi_NF_{N-1}\phi_{N-1}\cdots\phi_2\phi_1
      \bar\chi(t_0)\,\delta t_{N-1}                            \\
    &+\phi_N\phi_{N-1}F_{N-2}\phi_{N-2}\cdots\phi_2\phi_1
      \bar\chi(t_0)\,\delta t_{N-2}                            \\
    &+\cdots                                                     \\
    &+\phi_N\phi_{N-1}\cdots\phi_3F_2\phi_2\phi_1
      \bar\chi(t_0)\,\delta t_2                                \\
    &+\phi_N\phi_{N-1}\cdots\phi_3\phi_2F_1\phi_1
      \bar\chi(t_0)\,\delta t_1.
\end{aligned}
\label{eq:error-after-cancellation}
\end{equation}

\begin{equation}
    \Phi
    =\phi_N\phi_{N-1}\cdots\phi_2\phi_1.
    \label{eq:Phi}
\end{equation}

\begin{equation}
    \delta\boldsymbol t=
    \begin{bmatrix}
        \delta t_1\\
        \delta t_2\\
        \vdots\\
        \delta t_{N-1}\\
        \delta t_N
    \end{bmatrix}.
    \label{eq:delta-t-vector}
\end{equation}

\begin{equation}
\begin{aligned}
    \Gamma=\bigl[\,&
    \phi_N\phi_{N-1}\cdots\phi_3\phi_2F_1\phi_1\bar\chi(t_0),
    \\[2pt]&
    \phi_N\phi_{N-1}\cdots\phi_3F_2\phi_2\phi_1\bar\chi(t_0),
    \\[2pt]&
    \phi_N\phi_{N-1}\cdots\phi_4F_3\phi_3\phi_2\phi_1
      \bar\chi(t_0),
    \\[2pt]&
    \ldots,
    \\[2pt]&
    \phi_NF_{N-1}\phi_{N-1}\cdots\phi_2\phi_1\bar\chi(t_0),
    \\[2pt]&
    F_N\phi_N\phi_{N-1}\cdots\phi_2\phi_1\bar\chi(t_0)
    \,\bigr].
\end{aligned}
\label{eq:Gamma}
\end{equation}

\begin{equation}
\boxed{
    e(t_N)\doteq\Phi e(t_0)+\Gamma\,\delta\boldsymbol t
    }.
    \label{eq:final-linearization}
\end{equation}

\section*{Notes on corrections}

\begin{enumerate}
    \item The interval matrix is defined as
    $F_i=\widetilde A_{q_i}$, where $q_i$ is active on
    $[t_{i-1},t_i)$.  Writing $F_i=\widetilde A_{\sigma(t_i)}$ is ambiguous
    for a right-continuous switching signal because $\sigma(t_i)$ normally
    belongs to the following interval.

    \item The derivative in \eqref{eq:single-exponential-expansion} uses the
    augmented matrix $F_i$.  The physical matrix $A_{\sigma(t)}$ has the
    wrong dimension for the augmented exponential.

    \item The matrix order in \eqref{eq:ordered-product-expansion},
    \eqref{eq:error-after-cancellation}, and \eqref{eq:Gamma} has been
    corrected.  For example, the coefficient of $\delta t_{N-1}$ is
    \begin{equation*}
        \phi_NF_{N-1}\phi_{N-1}\cdots\phi_1\bar\chi(t_0),
    \end{equation*}
    not
    \begin{equation*}
        F_{N-1}\phi_N\phi_{N-1}\cdots\phi_1\bar\chi(t_0).
    \end{equation*}
    Matrices belonging to different modes do not generally commute.

    \item The columns of $\Gamma$ are written horizontally and in the same
    order as the entries of $\delta\boldsymbol t$.  This removes the
    dimensional ambiguity in the transposed stacked-vector notation.

    \item The symbol $\doteq$ means equality up to first order around
    \begin{equation*}
        e(t_0)=0,
        \qquad
        \delta\boldsymbol t=0.
    \end{equation*}
    Products such as $\delta t_i\delta t_j$ and
    $\delta t_i e(t_0)$ are omitted.

    \item Each $\delta t_i$ in this document is a change in an interval
    duration.  It is not yet an independent switching-instant offset.  If
    the initial and final cycle boundaries are fixed, the duration changes
    must also satisfy
    \begin{equation*}
        \delta t_1+\delta t_2+\cdots+\delta t_N=0.
    \end{equation*}
\end{enumerate}

\end{document}
